
\section{Introduction}

Embodied agents increasingly observe the world through streams rather than through clean symbolic states. A robot may receive a detector output that a cabinet is open, a relation estimate that a cup is inside the cabinet, an action log that a place operation succeeded, and a later visual observation that contradicts the earlier relation. These records may come from different sources, arrive at different times, and have different reliability. Some describe the current world state, while others arrive late and refer to an earlier event. Some evidence is missing because an object was occluded; some evidence is wrong because perception failed. A task executor nevertheless needs database-like answers: does the goal hold now, what would the system have known at a previous transaction time, which states changed during the episode, and which observations support or contradict an answer?

This is a data-management problem that is not isolated by current embodied benchmarks. BEHAVIOR introduced household activities with predicate-logic initial and goal conditions~\cite{srivastava2022behavior}, while BEHAVIOR-1K scales this setting to long-horizon mobile manipulation in a richer simulation environment~\cite{li2024behavior1k}. Habitat and Habitat 2.0 support embodied navigation and rearrangement evaluation~\cite{savva2019habitat,szot2021habitat}. AI2-THOR provides an interactive 3D environment for visual AI~\cite{kolve2017ai2thor}, and VirtualHome represents household activities as executable programs~\cite{puig2018virtualhome}. TEACh studies embodied agents that interact through dialogue~\cite{padmakumar2022teach}. OpenEQA evaluates open-vocabulary embodied question answering over episodic memory and active exploration~\cite{majumdar2024openeqa}, and ALFRED evaluates grounded instruction following for everyday tasks~\cite{shridhar2020alfred}. These benchmarks expose important task, language, perception, navigation, and interaction challenges. Their primary scoring target, however, is not whether a data system can maintain a temporal task-state view from imperfect public observation streams with valid-time semantics, transaction-time availability, uncertainty, provenance, and structured query workloads.

The distinction matters for database systems. In a task such as placing a cup into a cabinet and closing the door, the relevant maintained states may include \texttt{inside(cup,cabinet)}, \texttt{open(cabinet)}, and \texttt{grasped(robot,cup)}. The query \emph{was the cup inside the cabinet after the place action?} is not simply a nearest-frame retrieval query; it is a temporal state query over evidence. The query \emph{what changed between the beginning and end of the episode?} is not a single symbolic-state lookup; it is a change-set query over a maintained view. The query \emph{would the system have known this before the delayed observation arrived?} requires event time and arrival time to be separated, a distinction central to temporal data management~\cite{snodgrass1999developing,jensen1999temporal}. Since observations arrive continuously and may be processed out of order, the workload also inherits concerns from stream processing~\cite{babcock2002models}. Because task-state views must be updated after new evidence, corrections, and late arrivals, it is naturally related to incremental view maintenance~\cite{gupta1995maintenance,zhuge1995view}. Finally, because evidence can be incomplete, noisy, or contradictory, the benchmark connects to uncertain data and provenance-aware data management~\cite{aggarwal2009uncertain,green2007provenance,cheney2009provenance}.

We introduce \bench, a data-management benchmark for temporal task-state view maintenance over embodied observation streams. A \bench\ task instance provides public task specifications, StateObservation streams, queries, and metadata. It withholds hidden state timelines and hidden answer sets used by the evaluator. Systems under test may use any internal representation, including retrieval memory, temporal logs, SQL scans, rules, materialized views, learned memory, or hybrids. They are compared only through standardized QueryAnswer predictions under a shared evaluator. This public/hidden split is central to the benchmark: the oracle is derived from simulator truth and task specifications, not from any submitted system or reference implementation.

The reported \bench\ artifact contains 570 simulator-grounded episodes over 36 activities, 5,848 hard queries, 2.86M clean observations, and 2.80M hard mixed observations. The workload covers bitemporal state, provenance, change sets, uncertainty sets, preconditions, failure localization, and task-goal queries. Public-artifact validation, query-semantics audit, and answer-schema validation pass without errors. Experiments with recall-memory, temporal-log, and materialized temporal-state baselines show that the benchmark separates retrieval-style memory, log-based temporal lookup, and explicit temporal view maintenance: exact accuracy ranges from 25.2\% to 75.6\%, while structured, set-valued, and task-derived queries remain difficult.

We also provide \engine, a reference baseline engine for \bench. \engine\ maintains a \tsv\ with valid-time intervals, transaction-time revisions, support and contradict evidence, confidence, and task-derived views. It is useful for diagnosing what an explicit temporal view-maintenance engine can do under the benchmark, but it is not the benchmark oracle. This separation lets \bench\ evaluate \engine\ alongside simpler or future systems under the same public inputs, hidden answers, and scoring protocol.

This paper makes four contributions.
\begin{itemize}
    \item We formulate temporal task-state view maintenance over embodied observation streams as a data-management benchmark problem, separating it from perception, control, task completion, and memory-recall evaluation.
    \item We define the \bench\ artifact boundary, StateObservation schema, query workloads, perturbation regimes, QueryAnswer schema, and evaluator semantics for public-input/hidden-oracle evaluation.
    \item We describe a simulator-grounded generation pipeline that produces hidden state timelines, clean and perturbed public observation streams, query sets, hidden answer sets, validation reports, and reproducible benchmark manifests.
    \item We provide \engine\ as a reference baseline engine and an experimental protocol for comparing retrieval, log-based, and materialized temporal-state maintenance strategies under shared public inputs and hidden oracle answers.
\end{itemize}

The intended EAB contribution is therefore not a new perception model, planner, or robotics controller. It is the benchmark object: a standardized artifact, a diagnostic query workload, a public/hidden evaluation boundary, an evaluator, and an analysis protocol for systems that maintain temporal task-state views. This follows the tradition of database benchmark papers such as LDBC for graph workloads~\cite{erling2015ldbc,iosup2016ldbc}, TSM-Bench for time-series systems~\cite{khelifati2023tsmbench}, and FinBench for financial graph transaction workloads~\cite{qi2025finbench}. It also aligns with recent EAB work that contributes evaluation frameworks and diagnostic analyses, including SQLMorph~\cite{malekpour2026sqlmorph}, BEACON~\cite{najafi2025beacon}, and plan-based AQP analysis~\cite{mu2025aqp}.

